Interchange interventions test whether a network's internal states carry the variables an explanation assigns to
them. They are usually run one variable at a time, although causal abstraction requires agreement under
interventions on several variables at once. For alignments that map two variables to orthogonal subspaces at one
site, we show that a joint interchange edit equals a single edit applied to the result of the other. A joint test
can therefore reject an alignment that passes every single test only if both variables' single edits miss the natural
states they simulate by enough to push the output's margin below zero; a first-order version of this condition needs
only single queries. In transformers compiled with Tracr no alignment could fail only jointly. In small trained
transformers and MLPs, only \TrnPerfectJointFailWord{} of the \TrnPerfect{} alignments that pass every single test fail a
joint test, the first-order margin rule picks out all but one of them (post hoc), and at equal budgets joint tests
detected fewer failing alignments than single tests drawn from the queries in which the swapped variable changes
(\TrnContrastOneTwenty{} at 120 model calls). With the variables at different sites, joint tests did slightly better
(\MsContrast). In GPT-2 Small and GPT-2 Medium on indirect-object identification, under a protocol whose predictions
were hashed before any model was run, the prediction carried over from the small networks failed: joint tests detected
far more failing alignments than the same single tests (\ContrastS{} and \ContrastM{} at 120 calls). The position
subspace found by distributed alignment search also carries the names. Single queries whose donor keeps the position
reveal this, and every joint query exposes it; against single tests that also draw those queries, the advantage of
joint tests shrank to \ContrastAllS{} and \ContrastAllM. The \PerfectWordM{} GPT-2 Medium alignments that passed every
single query all failed jointly, and the margin rule flagged all of them.
